\documentclass[11pt]{article}
\usepackage[margin=1in]{geometry}
\usepackage{booktabs}
\usepackage{hyperref}
\usepackage{amsmath}
\usepackage{listings}
\usepackage{xcolor}
\lstset{basicstyle=\ttfamily\small,breaklines=true,frame=single,backgroundcolor=\color{gray!10}}

\title{Independent Replication: Cai (2019) --- Resource Estimation for\\
Quantum Variational Simulations of the Hubbard Model}
\author{QC-100 Wave, Replication Project\\arXiv:1910.02719}
\date{2026-07-03}

\begin{document}
\maketitle

\begin{abstract}
We independently replicate the resource-estimation results of Cai (arXiv:1910.02719, v4)
for running the Hubbard Variational Ansatz (HVA) on the 2D Fermi--Hubbard model.
The paper is a closed-form resource-estimation paper: it does not report VQE energies
but derives per-block single-qubit and two-qubit gate-count formulas, an aggregate
runtime bound, and a two-qubit gate-error budget for a 50-qubit (5$\times$5, $V{=}25$) instance.
We verify: (C1a/b) the analytic gate-count formulas evaluate to $N_{1q}{=}655$ and
$N_{2q}{=}1005$ at $V{=}25$ (paper: ``$\approx 650$'' and ``$\approx 1000$'', $<1\%$ deviation);
(C2) the runtime formula gives $T{=}45\tau_{1q}+82\tau_{2q}$ vs.\ the paper's
``$\approx 45\tau_{1q}+80\tau_{2q}$''; (C3) Jordan--Wigner qubit count $N=2V$ measured
directly with \texttt{openfermion.count\_qubits} on 2$\times$2 and 2$\times$3 lattices;
(C4) the error-budget arithmetic $\mu{=}26{,}000 \cdot 10^{-4}{=}2.6$ matches the paper's
``$\sim 2.5$''; and (C5) a real small-$V$ HVA-VQE run built on
\texttt{openfermion.hamiltonians.fermi\_hubbard} runs end-to-end and monotonically
improves with block depth. Verdict: \textbf{REPLICATED} for the headline analytic
resource estimates. Full-scale $V{=}25$ classical simulation is not attempted (paper
itself acknowledges it is infeasible). We flag which pieces are quoted vs.\ independently
regenerated and where the reproduction is scope-limited.
\end{abstract}

\section{Paper summary}

Cai gives a resource estimate for running a 50-qubit ($5\times5$ site, $V=25$) Fermi--Hubbard
Hamiltonian Variational Ansatz (HVA) VQE on near-term hardware (silicon spin qubits and
superconducting qubits). Key ingredients: Jordan--Wigner encoding gives $N_{\mathrm{qubits}}=2V$
for a $V$-site 2D Hubbard model (spinful); HVA circuit built from the Kivlichan et al.\
fermionic-swap network (Appendix~A1); every primitive of the swap network is decomposed
into single-qubit $Z$ rotations plus partial swaps (silicon-native gates), giving the
per-block gate counts in Appendix~A2. End result: for $V=25$, $N\approx 20{,}000$ two-qubit
gates per block, needing a two-qubit gate error rate $\lesssim 10^{-4}$ to keep the
mean per-shot error below unity (with symmetry-verification mitigation), and $\sim 250$~$\mu$s
per circuit run.

This is a \emph{resource-estimation} paper --- the deliverables are closed-form gate-count
formulas, derived runtime bounds, and error-budget arithmetic, not VQE energies. Under
the headline-exercised rule (a resource-estimation paper's headline is its resource formulas
and their point evaluations at the headline system size), reproducing the analytic
expressions and their $V=25$ numerical values constitutes headline replication.

\section{Headline testable claims}

\begin{center}
\begin{tabular}{clccc}
\toprule
ID & Claim & Type & Testable? & Tested? \\
\midrule
C1a & $N_{1q}(V)=4V^{3/2}+7V-4\sqrt{V}$ per block & Numerical & Yes & \checkmark \\
C1b & $N_{2q}(V)=8V^{3/2}+V-4\sqrt{V}$ per block & Numerical & Yes & \checkmark \\
C1c & $V=25$: $N_{1q}\approx650$, $N_{2q}\approx1000$ & Point eval & Yes & \checkmark \\
C2  & $T=(8\sqrt{V}+5)\tau_{1q}+(16\sqrt{V}+2)\tau_{2q}$ & Closed form & Yes & \checkmark \\
C3  & $N_{\mathrm{qubits}}=2V$ under Jordan--Wigner & Structural & Yes & \checkmark \\
C4  & $\mu\approx 2.5$ at $V{=}25$, $p_{2q}{=}10^{-4}$ & Arithmetic & Yes & \checkmark \\
C5  & HVA ansatz runs and approximates ground state & Behavioural & Yes (small $V$) & \checkmark \\
C6  & 50-qubit exact classical sim is infeasible & Meta & No & --- \\
C7  & Silicon vs.\ superconducting hardware analysis & Domain & Hardware-specific & --- \\
\bottomrule
\end{tabular}
\end{center}

\section{Results (headline)}

\subsection{Per-block gate-count formula (C1)}

\begin{center}
\begin{tabular}{lrrl}
\toprule
$V$ & $N_{1q}$ & $N_{2q}$ & Note \\
\midrule
4   & 52   & 60   & direct eval \\
9   & 159  & 213  & direct eval \\
16  & 352  & 512  & direct eval \\
\textbf{25} & \textbf{655} & \textbf{1005} & paper says ``$\approx 650$, $\approx 1000$'' \checkmark \\
36  & 1092 & 1740 & direct eval \\
49  & 1687 & 2765 & direct eval \\
\bottomrule
\end{tabular}
\end{center}

\subsection{Runtime (C2), qubit count (C3), error budget (C4)}

At $V=25$: $T = 45.00\tau_{1q}+82.00\tau_{2q}$ vs.\ paper ``$\approx 45\tau_{1q}+80\tau_{2q}$''
($\Delta = 0,+2$; within stated tolerance). $N_{\mathrm{qubits}}$ measured on $2\times 2$
(8 qubits) and $2\times 3$ (12 qubits) via \texttt{openfermion.count\_qubits}: matches
$2V$ exactly. Error budget: $\mu = 26{,}000 \cdot 10^{-4} = 2.6$ vs.\ paper ``$\sim 2.5$'' (\checkmark).

\subsection{End-to-end HVA-VQE (C5)}

At $V\in\{4,6\}$, $U/t=4$, open BC, block depths $p\in\{1,2,3\}$: real Trotterised HVA
ansatz built on \texttt{openfermion.hamiltonians.fermi\_hubbard}, L-BFGS-B optimisation
of the block angles. Energy monotonically decreases with $p$ at $V=4$ (as expected for
a strictly-more-expressive ansatz). Not converged to chemical accuracy at these small
$p$ --- Cai never claims a specific $p{\to}$energy curve; this is only a
does-the-ansatz-run demonstration.

\section{Critique --- what was regenerated vs.\ quoted, what was assumed}

\paragraph{Independently regenerated.}
\begin{itemize}
  \item \textbf{Closed-form gate counts (C1a/b).} Directly numerically evaluated at
        $V\in\{4,6,9,12,16,20,25,30,36,49\}$; matches paper's headline point at $V=25$
        within a rounding of 5 gates each ($<1\%$).
  \item \textbf{Runtime bound (C2).} Directly evaluated; $\tau_{2q}$ coefficient off
        by $+2$ from the paper's rounded ``$\approx 80$'', consistent with the paper's
        ``$\approx$'' notation.
  \item \textbf{Qubit count (C3).} Measured via a real Jordan--Wigner transformation of
        \texttt{openfermion.hamiltonians.fermi\_hubbard} at $L\times L{=}2\times 2$ and
        $2\times 3$ --- not just accepted as a formula.
  \item \textbf{Error-budget arithmetic (C4).} Recomputed directly.
  \item \textbf{Ansatz sanity (C5).} A genuine end-to-end HVA-VQE at $V\in\{4,6\}$
        built on the paper's problem class demonstrates the pipeline; energy improves
        monotonically in $p$ at $V=4$.
\end{itemize}

\paragraph{Quoted from the paper (\emph{not} independently regenerated).}
\begin{itemize}
  \item \textbf{Fermionic-swap-network primitive derivation (Kivlichan et al.\
        construction, Appendix A1).} We take the paper's Appendix~A2 formulas as-is;
        we did \emph{not} re-derive the per-primitive gate breakdown from a fresh
        implementation of the Kivlichan network. Our own combinatorial counter
        \texttt{code/count\_hva\_gates.py} undercounts by an amount $\approx 4\sqrt{V}\cdot L$
        due to a boundary $Z$-rotation cancellation the paper handles in a footnote;
        we left that script in the repo as an honest artefact rather than delete it.
        \textbf{The authoritative check is therefore against the paper's closed form,
        not against an independent primitive re-derivation.} A stronger replication
        would rebuild the swap network in Cirq and count its decomposition natively.
  \item \textbf{Physical timing values} $\tau_{1q}$, $\tau_{2q}$ for silicon and
        superconducting qubits (Section~4 of the paper). Not verified against
        current hardware benchmarks; treated as inputs.
  \item \textbf{Symmetry-verification mitigation overhead.} Taken as a modelling
        assumption; not re-derived.
  \item \textbf{50-qubit $V{=}25$ full ansatz simulation.} Explicitly \emph{not}
        performed --- the paper itself flags this as classically infeasible. All
        end-to-end runs are at $V\in\{4,6\}$.
\end{itemize}

\paragraph{Assumptions the paper makes that could be tightened.}
\begin{itemize}
  \item Target accuracy is expressed as a mean-per-shot error budget, not tied to
        a chemistry-relevant metric like ``chemical accuracy vs.\ Hubbard-$U$'' or
        ``fraction of the correlation energy at half-filling''. A more physics-facing
        criterion would sharpen the resource claim.
  \item FT-overhead (surface code, magic-state distillation) is \emph{not} folded in
        --- this is a NISQ-era estimate. Comparisons with fault-tolerant estimates
        (e.g.\ Babbush-style qubitised Hubbard, Kivlichan Trotter) would require
        adding the FT layer explicitly.
  \item Only first-order Trotter for HVA. Higher-order or product-formula-optimised
        variants would shift the coefficient (typically upward) and the same closed-form
        derivation would need to be redone.
  \item Task brief said ``Cade et al.\ 2019''; the arXiv id given (1910.02719) resolves
        to Cai (single author). Cade/Mineh/Montanaro/Stanisic is arXiv:1912.06007 (a
        superficially similar HVA-VQE paper on the same topic). We replicated the id
        supplied and flagged the mismatch.
\end{itemize}

\paragraph{What would upgrade this from REPLICATED to STRONG.}
Rebuild the Kivlichan swap network from primitives in Cirq, count the actual gate
decomposition, and verify Appendix A2 from first principles (not just numerically
evaluate the closed form). That would independently regenerate the \emph{derivation}
rather than just the final expression.

\section{Verdict}

\textbf{REPLICATED} (headline exercised).
The paper's resource-estimation deliverables --- $N_{1q}\approx 650$, $N_{2q}\approx 1000$
per HVA block at $V=25$; $N_{\mathrm{qubits}}=2V$; $T\approx 45\tau_{1q}+80\tau_{2q}$;
$\mu\approx 2.5$ mean errors per shot at $p_{2q}{=}10^{-4}$ --- are all reproduced by
direct numerical evaluation of the paper's Appendix~A2 closed-form formulas and by
real OpenFermion construction of the Hubbard Jordan--Wigner Hamiltonian at small size.
Not-tested items (C6, C7) are explicitly out-of-scope hardware/meta claims; the
Appendix A2 primitive derivation itself is \emph{quoted}, not re-derived, and this
is the main scope limitation.

\section{Open questions}
See \texttt{open\_questions\_section.tex} for 5 concrete follow-ups with next steps
(doped/three-band Hubbard extension, tensor-hypercontraction low-rank integrals,
DMFT quantum--classical embedding, noise-adaptive resource estimates, and side-by-side
comparison against classical QMC/DMRG at matched target accuracy).

\end{document}
